\section{Security Analysis}

\label{sec:security}
% ============================================================================

\begin{theorem}[Threshold Security]
The Hanzo Threshold Signing system with $(t, n) = (3, 5)$ provides security
under the following adversary model: an adversary who compromises up to $t-1 = 2$
signing nodes cannot forge signatures, recover the private key, or distinguish
partial signatures from random values.
\end{theorem}

\begin{remark}
The identifiable abort property of CGGMP21 ensures that a compromised node
cannot cause signing to fail silently. If a node submits invalid partial
signatures, the coordinator detects the fault, identifies the malicious node,
and re-runs the protocol with the remaining honest nodes (provided $\geq t$
honest nodes remain).
\end{remark}

\subsection{Deployment Topology}

The 5 signing nodes are deployed across 3 geographic regions and 2 cloud
providers:
\begin{itemize}[leftmargin=1.1em]
    \item 2 nodes in US (DOKS lux-k8s, separate availability zones).
    \item 2 nodes in EU (DOKS lux-k8s, separate availability zones).
    \item 1 node in APAC (DOKS lux-k8s).
\end{itemize}

An adversary must compromise nodes in at least 2 geographic regions to reach
the threshold.

% ============================================================================
